\begin{lstlisting}[
language=Python,
caption={Simulation of a spiking neuron group using Brian2.},
label={lst:brian_neuron_sim},
basicstyle=\ttfamily\footnotesize,
breaklines=true,
breakatwhitespace=false,
columns=fullflexible,
frame=single,
captionpos=b,
numbers=left,
numberstyle=\tiny\color{gray},
keywordstyle=\color{blue},
commentstyle=\color{gray},
stringstyle=\color{orange},
showstringspaces=false
]
from brian2 import *

# Simulation duration
duration = 1*second

# Input: Poisson spike train with constant firing rate
rate_input = 100*Hz

# Leaky Integrate-and-Fire (LIF) neuron parameters
tau = 10*ms              # Membrane time constant
u_rest = 0*mV            # Resting potential
u_reset = u_rest         # Reset potential after spike
u_threshold = 20*mV      # Spike threshold

# Neuron model equation
eqs = '''
du/dt = (u_rest - u) / tau : volt (unless refractory)
'''

# Create a group of N neurons
N = 10
neurons = NeuronGroup(
    N, model=eqs,
    threshold='u > u_threshold',
    reset='u = u_reset',
    refractory=100*ms,
    method='exact'
)

# Initialize membrane potentials randomly
neurons.u = 'u_rest + rand() * (u_threshold - u_rest)'

# Poisson input group for rate encoding
input_group = PoissonGroup(N, rates=rate_input)

# Synaptic connections: one-to-one, with fixed PSP
synapses = Synapses(input_group, neurons, on_pre='u += 10*mV')
synapses.connect(j='i')

# Record spikes and membrane potentials
spike_monitor = SpikeMonitor(neurons)
voltage_monitor = StateMonitor(neurons, 'u', record=True)

# Run simulation
run(duration)

# Plotting results
plt.figure(figsize=(10, 4))

# Spike raster plot
plt.subplot(121)
plt.title("Spike Raster")
plt.plot(spike_monitor.t/ms, spike_monitor.i, '.k')
plt.xlabel('Time (ms)')
plt.ylabel('Neuron index')

# Membrane potential trace of first neuron
plt.subplot(122)
plt.title("Voltage Trace (Neuron 0)")
plt.plot(voltage_monitor.t/ms, voltage_monitor.u[0]/mV)
plt.xlabel('Time (ms)')
plt.ylabel('Membrane potential (mV)')

plt.tight_layout()
plt.show()
\end{lstlisting}
